\subsection{The complete first-level range in residue characteristic three}

\begin{lemma}[B-RD9: a variable-ramification Hesse family]\label{lem:B-RD9}
\status{Proved.}
Let $K$ have algebraically closed residue field of characteristic three,
absolute ramification index $e$, and a uniformizer $\varpi$. Suppose
$\mu_3\subset K$. Choose an integer $r$ with
$0<2r<e$ and $3\nmid e+r$, and put $t=1+\varpi^r$. The Hesse
Weierstrass curve of Theorem~\ref{thm:B-RD8} has full rational
$3$-torsion and semistability defect
\[
 m=\frac{12}{\gcd(4,e+r)}.
\]
\end{lemma}
\begin{proof}
The torsion formulas in Theorem~\ref{thm:B-RD8} are identities in $t$,
so apply unchanged. Write $w=\varpi^r$ and first scale the $2$-division
polynomial by $x=3X$. After translating $X=w^2+Z$, the resulting cubic is
\[
 Z^3+3w^2Z^2-(12w^3+18w^2+36w+27)Z
       +12w^4-36w^3-117w^2-108w-54.
\]
Its coefficients, from constant through quadratic, have valuations
$e+4r,e+3r,e+2r$. The inequalities $2r<e$ make the displayed leading
terms unique. Its Newton polygon has the single slope
$-(e+4r)/3$. Because $3\nmid e+r$, adjoining a root $z$ gives a
completely ramified cubic extension $K_1/K$. Put $x_0=3(w^2+z)$.
Then
\[
 v_K(x_0)=e+2r,\qquad v_K(3x_0^2+a)=3(e+r).
\]
The derivative computation giving the second equality has a unique
lowest term, of valuation $e+3r$, before multiplying by $3^2$.

Choose $u$ with $u^4=3x_0^2+a$. Since the residue field is algebraically
closed and $4$ is invertible, adjoining this root has ramification
degree $4/\gcd(4,e+r)$ over $K_1$. The equation obtained from
$x=x_0+u^2X$, $y=u^3Y$ is
\[
 Y^2=X^3+(3x_0/u^2)X^2+X,
\]
and the middle coefficient has valuation $(e+r)/2>0$. Thus it is good.
Conversely $v_K(\Delta)=9(e+r)$, which forces the tame factor
$4/\gcd(4,e+r)$ in any good-reduction ramification degree. The
$2$-division cubic forces the additional factor three, exactly as in
Theorem~\ref{thm:B-RD8}. This proves minimality.
\end{proof}

\begin{theorem}[B-RD10: all first-level defects at $p=3$]\label{thm:B-RD10}
\status{Proved; independently audited by the parent agent.}
Allowing arbitrary perfect residue fields, a potentially good additive
elliptic curve with full rational $3$-torsion exists at every positive
even absolute ramification index. The precise possible defects are
\[
\begin{array}{c|c}
\text{defect}&\text{absolute ramification indices}\\\hline
2& e\in2\mathbf Z,\ e\geq4\cr
3& e\in2\mathbf Z,\ e\geq4\cr
4& e\in2\mathbf Z,\ e\geq2\cr
6& e\in2\mathbf Z,\ e\geq2\cr
12& e\in2\mathbf Z,\ e\geq2.
\end{array}
\]
Over the particular field $\mathbf Q_3(\zeta_3)$, Theorem~\ref{thm:B-RD2}
already gives an example of defect six.
\end{theorem}
\begin{proof}
The special fibre is supersingular by Proposition~\ref{prop:B-RD3}.
Over an algebraically closed field of characteristic three it is
isomorphic to $y^2=x^3-x$. Its automorphisms are
\[
 (x,y)\longmapsto(u^2x+r,u^3y),\qquad
 u^4=1,\quad r^3-r=0.
\]
Thus its automorphism group has order twelve, is $C_3\rtimes C_4$,
and the only possible nontrivial subgroup orders are $2,3,4,6,12$.
The inertia image embeds in this group.

Over a good-reduction extension of ramification degree $m$, all eight
nonzero $3$-torsion parameters have positive integral valuations.
Their sum is $v_L(3)=me$, by the constant term of $[3](T)/T$.
Consequently $me\geq8$, which rules out $m=2,3$ at $e=2$.
The Weil pairing gives $\mu_3\subset K$, hence $e$ is even.
These are all the necessary restrictions in the table.

The tame cases $m=2,4$ are Theorem~\ref{thm:B-RD7}.
We explain the remaining constructions explicitly. For any desired
even $e$, a convenient base with algebraically closed residue field is
\[
 K_e=\operatorname{Frac}\left(
 W(\overline{\mathbf F}_3)[\varpi]/
 (\varpi^e+3\varpi^{e/2}+3)\right).
\]
The polynomial is Eisenstein, and $1+\varpi^{e/2}$ is a primitive
cube root of unity. In all the cases below the variable-ramification
family is taken over $K_e$.

For defect twelve and $e\geq4$, use $r=1$ unless $3\mid e+1$, in
which case use $r=3$. The second case first occurs at $e=8$, so always
$0<2r<e$. The integer $e+r$ is odd and prime to three. The missing
case $e=2$ is Theorem~\ref{thm:B-RD8}.

For defect six with $e\equiv0\pmod4$, use $r=2$, replacing it by
$r=6$ when $3\mid e+2$. These choices work for every such $e\geq8$.
For $e\equiv2\pmod4$, use $r=4$, replacing it by $r=8$ when
$3\mid e+4$. These choices work except at $e=2,6,14$.
In each applicable case, $e+r\equiv2\pmod4$ and $3\nmid e+r$,
so Lemma~\ref{lem:B-RD9} applies. The case $e=2$ is
Theorem~\ref{thm:B-RD2}. The case $e=4$ is obtained from the
order-twelve example by the tame quadratic extension: its inertia
image then has order six. The case $e=14$ follows from the order-six
example at $e=2$ by a tame degree-seven extension.
For $e=6$, start with the order-six example over a base of algebraically
closed residue field, and adjoin a cube root of a uniformizer.
Its good-reduction cubic subfield is obtained by adjoining a cube root
of $-1/4$, whose Kummer class has valuation zero; the uniformizer class
has valuation one. These cyclic cubic extensions are distinct and
linearly disjoint. Thus the order-six inertia image remains unchanged.
All finite extensions here are totally ramified.

For defect three with $e\equiv0\pmod4$, use $r=4$, replacing it by
$r=8$ when $3\mid e+4$. These choices work for every such $e\geq12$.
For $e\equiv2\pmod4$ and $e\geq6$, use $r=2$, replacing it by
$r=6$ when $3\mid e+2$; only $e=10$ is exceptional.
The examples at $e=4,8$ come from the order-six example at $e=2$
by tame extensions of degree two and four: they remove its order-two
inertia and retain the order-three subgroup.
For $e=10$, work over $K_{10}$ and put $t=1+\varpi^2$.
In the $2$-division cubic scaled by $x=3X$, translate
$X=\varpi^4+\varpi^6+Z$. Its constant, linear and quadratic
coefficients have valuations $22,16,14$, respectively. Its Newton
polygon therefore has the single slope $-22/3$, forcing a ramified
cubic extension. Here $v_K(\Delta)=108$, and the good-model
construction has $v_K(u^4)=36$; its fourth root requires no additional
ramification over the cubic extension. Hence the defect is exactly
three. The exact certificate \texttt{verify\_p3\_ranges.py} verifies these
coefficient valuations by reducing polynomials modulo
$\varpi^{10}+3\varpi^5+3$, and separately audits all the index recipes
for every even $e\leq10000$.

For the assertions about tame base change, one can work throughout
over completed maximal unramified bases. The order-six example has
cyclic inertia $C_6$. The order-twelve example has the full group
$C_3\rtimes C_4$, whose tame quotient is $C_4$; a tame quadratic
extension leaves the order-six subgroup. These observations also
verify that the stated defects are not merely upper bounds.
\end{proof}

\begin{corollary}[B-RD11: the cyclotomic field]\label{cor:B-RD11}
\status{Proved.}
For $p\geq5$, no elliptic curve over $\mathbf Q_p(\zeta_p)$ can have
full rational $p$-torsion and additive potentially good reduction.
For $p=3$, such a curve exists.
\end{corollary}
\begin{proof}
For $p\geq7$, the necessary inequality in
Theorem~\ref{thm:B-RD7} gives
$e\geq(p^2-1)/6>p-1$, contradicting $e(\mathbf Q_p(\zeta_p))=p-1$.
For $p=5$, the only possible defect is six. We justify that the good
extension may be chosen totally ramified, so its residue field really
is $\mathbf F_5$. Choose a minimal integral short Weierstrass equation
$y^2=x^3+ax+b$, and put $\delta=v_K(\Delta)$. Since the reduction is
potentially good, $v_K(j)\geq0$, giving $3v_K(a)\geq\delta$.
The discriminant identity then gives $2v_K(b)\geq\delta$ as well.
Put $m'=12/\gcd(12,\delta)$. Over the totally ramified extension
$K(\varpi^{1/m'})$, scaling by $u=\varpi^{\delta/12}$ gives integral
coefficients and unit discriminant. Conversely the discriminant
forces $m'$ to divide every good-reduction ramification degree.
Thus $m'=6$, and this construction gives the required degree-six
extension with unchanged residue field. Equality now holds in the
valuation bound: over this good-reduction extension every nonzero
$5$-torsion formal parameter has valuation one. Dividing those
parameters by a uniformizer and reducing gives an injective
$\mathbf F_5$-linear map
\[
 E[5](K)\longrightarrow k_L=k_K=\mathbf F_5.
\]
Additivity follows from the linear term of the formal group law;
injectivity follows because every nonzero parameter has valuation
one. Its source would have dimension two and its target dimension
one, a contradiction. The claim for $p=3$ is
Theorem~\ref{thm:B-RD2}.
\end{proof}
